\cleardoublepage
\chapter*{Abstract}
\addcontentsline{toc}{chapter}{Abstract}

\textbf{Title:} Semi-supervised model of equivariant convolutional neural networks for segmentation and classification of histopathological images.\footnote{Bachelor thesis.}

\textbf{Author:} \Author.\footnote{Faculty of Physical-Mechanical Engineering. School of Systems Engineering and Informatics. Advisor: \DirectorName, \DirectorDegree.}

\textbf{Keywords:} equivariant convolutional neural networks, histopathological images, digital image processing, semi-supervised learning.

\textbf{Description:}

\begingroup\small

% TODO: translate the approved final Spanish abstract; do not maintain two
% independently edited scientific claims.
Digital pathology represents a paradigm shift in the traditional approach to medical diagnosis. Within contemporary medicine, ongoing technological advances are constantly reshaping healthcare practices. Therefore, the development of digital pathology is a transformative achievement that allows for the introduction of a combination of cutting-edge electronic devices, such as histology slide scanners, data analytics, and clinical expertise that can significantly transform the way patient care is delivered.

Among the applications of digital pathology, one of the most popular and interesting has been that regarding cancer. This is because it is one of the leading causes of death globally. The burden of cancer will increase by approximately 60\% over the next two decades, further affecting healthcare systems, individuals, and communities.

Given this, a new approach is required, such as the use of artificial intelligence (AI) models, which optimize and accelerate current methods. Although machine learning projects have been developed within the field of digital pathology, the limited availability of annotated and segmented images hinders deep learning processes, since performing these annotation processes requires a great deal of time and effort from pathologists. Until now, traditional deep learning methods have addressed these problems by using data augmentation, which consists of generating new data from existing data, but this project aims to show the use of equivariant representations in deep learning models, which would allow a stable and effective representation of images in digital pathology.

For this reason, a semi-supervised deep learning model of equivariant neural networks is proposed to optimize preprocessing time and resources, while reducing the limitations that pathologists face when implementing automatic analysis projects.
\endgroup
